% --- Vorwort der Veranstalter*innen ---
\clearpage
\phantomsection
\addcontentsline{toc}{part}{Vorwort}
\section*{Vorwort}

Liebe WissKI-Community,

mit dem vorliegenden \emph{Book of Abstracts} möchten wir Ihnen einen kompakten Überblick über die Beiträge geben, die im Rahmen unseres diesjährigen Treffens an der Friedrich-Alexander-Universität Erlangen-Nürnberg vorgestellt werden. Das Anwender*innentreffen ist mittlerweile fester Bestandteil im Kalender der WissKI-Community und versteht sich als Forum für den fachlichen Austausch zwischen Forschenden, Ent\-wick\-ler*in\-nen, Datenkurator*innen und allen, die WissKI in ihren digitalen Editions-, Sammlungs- oder Forschungsprojekten einsetzen.

In diesem Jahr haben uns Einreichungen aus einem breiten Spektrum geisteswissenschaftlicher Disziplinen erreicht. Sie zeigen eindrücklich, wie vielfältig die Anwendungskontexte sind, in denen \emph{Wissenschaftliche Kommunikationsinfrastrukturen} heute zum Einsatz kommen -- von der musealen Sammlungsdokumentation über digitale Editionen und Forschungsdatenmanagement bis hin zur Anbindung an Normdaten- und Linked-Open-Data-Ökosysteme. Die Beiträge sind in drei Kategorien gegliedert: \textbf{Präsentationen} berichten über abgeschlossene oder weit fortgeschrittene Arbeiten und Erkenntnisse; \textbf{Workshops} laden zum praxisnahen Mitarbeiten an konkreten Werkzeugen und Methoden ein; und im Format \textbf{Bring your WissKI} stellen Projekte ihre laufenden Arbeiten vor und bringen sie zur gemeinsamen Diskussion.

Wir danken allen Autor*innen für ihre Beiträge, dem Programmkomitee für die sorgfältige Begutachtung und der WissKI-Community für ihr fortwährendes Engagement. Unser besonderer Dank gilt den Förderern und institutionellen Partnern, ohne deren Unterstützung das Treffen nicht möglich wäre.

Wir wünschen Ihnen eine anregende Lektüre, fruchtbare Gespräche und ein erfolgreiches Anwender*innentreffen 2026.

\vspace{1cm}

\noindent\textit{Die Veranstalter*innen}\\
Nürnberg, im Oktober 2026
